\section{Search policy: from plausible continuation to obstruction removal}

At each point the system identifies a minimal or approximate set of unresolved blockers intersecting viable proof routes: a proof-theoretic analogue of an epistemic cut set. Candidate actions are selected for their expected ability to remove blockers or split ambiguity. High-value actions include proving a reusable bottleneck lemma, finding a counterexample that kills a family of conjectures, changing coordinates, importing a structurally analogous invariant, weakening an overstrong conjecture, or locating a known parent theorem.

This architecture aligns with recent successful discovery systems. FunSearch couples a pretrained model to systematic evaluation and evolutionary selection, producing new constructions and algorithms \cite{romera2024}. AlphaEvolve expands this evaluator-guided search pattern to scientific and algorithmic discovery \cite{novikov2025}. AlphaProof treats formal proof search as sequential decision making inside Lean \cite{hubert2025}. Large-scale AlphaEvolve experiments report both rediscovery of known best solutions and improvements on several mathematical problems, highlighting why discovery and novelty classification must be separated \cite{georgiev2025}. The common lesson is that generation becomes more reliable when downstream state transitions are controlled by evaluators.
